\documentclass{article}

\usepackage[utf8]{inputenc}     % pdfLaTeX
\usepackage[T1]{fontenc}
\usepackage[magyar]{babel}

% Set page size and margins
% Replace `letterpaper' with `a4paper' for UK/EU standard size
\usepackage[a4paper,top=2cm,bottom=2cm,left=3cm,right=3cm,marginparwidth=1.75cm]{geometry}
%for arrows
\usepackage{textcomp}

\title{Földrajzi nevek helyesírása}
\author{Lengyel Zoltán Péter \thanks{Csukás Ibolya tanítása alapján}}
\date{Leguóbb frissítve: 2026. Március 11.}

\begin{document}

\maketitle

\tableofcontents
\newpage

\section{Egybeírt földrajzi nevek}
Egybeírjuk a következő utótagú országneveket és nagy tájegységek nevét:
\begin{itemize}
    \item ország (pl. Magyarország)
    \item föld (pl. Thaiföld)
    \item alföld (pl. Németalföld)
    \item part (pl. Elefántcsontpart)
    \item alja (pl. Alpokalja)
\end{itemize} 

\subsection{-i képzős alakok}
Az -i képzős alakot kisbetűvel írjuk
\begin{itemize}
    \item pl. magyarországi, thaiföldi, németalföldi, elefántcsontparti, alpokaljai
\end{itemize}

\subsubsection{i-re végződő alap alakok}
i-re végződő alap alakok esetén csak egy i-t írunk
\begin{itemize}
    \item Helsinki\textrightarrow helsinki
    \item Tamási\textrightarrow tamási
\end{itemize}

\section{Kötőjellel írt földrajzi nevek}
\subsection{Ha az utótag földrajzi köznév}
Ha az utótag földrajzi köznév, az előtagot nagybetűvel, az utótagot kisbetűvel írjuk
\begin{itemize}
    \item óceán 
    \item tenger 
    \item tó 
    \item folyó
    \item folyam
    \item patak
    \item hegy
    \item (közép)hegység
    \item domb
    \item dombság
    \item völgy
    \item hágó
    \item dűlő
    \item fennsík
    \item sivatag
    \item sziget(ek)
    \item félsziget
    \item öböl
    \item stb.
    \item pl. Velencei-tó, Sólyom-sziget, Szabadság-hegy, Dunántúli-középhegység, Csendes-óceán stb.
\end{itemize}

\subsection{-i képzős alakok}
\subsubsection{Ha az előtag is köznév}
Ha az előtag is köznév, mindkét tagot kisbetűvel írjuk
\begin{itemize}
    \item Sólyom-sziget\textrightarrow sólyom-szigeti
    \item Szabadság-hegy\textrightarrow  szabadság-hegyi
    \item Velencei-tó\textrightarrow  velencei-tavi
\end{itemize}

\subsubsection{Ha az előtag tulajdonnév}
Ha az előtag tulajdonnév, az első tagot nagybetűvel, a második tagot kisbetűvel írjuk.
\begin{itemize}
    \item János-hegy\textrightarrow János-hegyi
    \item Duna-part\textrightarrow Duna-parti
    \item Balaton-felvidék\textrightarrow Balaton-felvidéki
    \item Margit-sziget\textrightarrow Margit-szigeti \footnote{itt Margit-sziget mint a sziget, amennyiben a városrészről van szó, Margitsziget}
\end{itemize}

\subsubsection{Ismeretlen eredetű előtag}
Ismeretlen eredetű előtag esetében az első tagot nagybetűvel, a második tagot kisbetűvel írjuk
\begin{itemize}
    \item Kab-hegy\textrightarrow Kab-hegyi
    \item Kaszpi-tenger\textrightarrow Kaszpi-tengeri
\end{itemize}

\subsection{Ha az utótag tulajdonnév}
Ha az utótag tulajdonnév, az előtag vagy melléknév, vagy égtáji elnevezés, akkor mindkét tagot nagybetűvel írjuk
\begin{itemize}
    \item pl. Fekete-Kőrös, Magas-Tátra, Holt-Tisza, Új-Zéland, Közép-Európa, Délnyugat-Magyarország, Nyugat-Dunántúl, Délkelet-Ázsia
\end{itemize}

\subsubsection{-i képzős alakok}
-i képzős alakban mindkét tagot kisbetűvel írjuk.
\begin{itemize}
    \item Közép-Európa\textrightarrow közép-európai
    \item Magas-Tátra\textrightarrow magas-tátrai
    \item Holt-Tisza\textrightarrow holt-tiszai
    \item Új-Zéland\textrightarrow új-zélandi
    \item Nyugat-Dunántúl\textrightarrow nyugat-dunántúli
\end{itemize}

\section{Különírt földrajzi nevek}
\subsection{Mai és történelmi államnevek}
Mai és történelmi államok neveiben mindkét tagot nagybetűvel írjuk
\begin{itemize}
    \item pl. Római Birodalom, Kijevi Nagyfejedelemség, Egyesült Királyság, Amerikai Egyesült Államok
\end{itemize}

\subsubsection{-i képzős alakok}
-i képzős alakban minden tagot kisbetűvel írunk
\begin{itemize}
    \item pl. római birodalmi, kijevi nagyfejedelemségi, egyesült királysági, amerikai egyesült államokbeli
\end{itemize}

\subsection{Közigazgatási utótag}
Közigazgatási utótagú földrajzi nevekben az előtagot nagybetűvel, az utótagot pedig kisbetűvel írjuk

\begin{itemize}
    \item (vár)megye, régió, járás
    \begin{itemize}
        \item   pl. Győr-Moson-Sopron vármegye
    \end{itemize}
    \item út, utca, tér, köz, körút, körtér, fasor, lépcső
    \begin{itemize}
        \item pl. Széchenyi utca, Batthyány tér, Gorkij fasor, Szent István Körút
    \end{itemize}
    \item híd
    \begin{itemize}
        \item pl. Margit híd, Szabadság híd \footnote{Kivétel a Lánchíd}
    \end{itemize}
    \item köze, mente, vidéke, környéke, melléke
    \begin{itemize}
        \item Duna-Tisza köze, Szekszárd vidéke, Dráva mente, Pest környéke
    \end{itemize}
\end{itemize}

\subsubsection{-i képzős alakok}
-i képzős alakokban az előtagot nagybetűvel, az utótagot kisbetűvel írjuk
\begin{itemize}
    \item pl. Győr-Moson-Sopron vármegyei, Szent István körúti, Margit hídi \footnote{Lánchíd\textrightarrow lánchídi}, Duna-Tisza közei, Dráva menti
\end{itemize}

\end{document}